% =====================================================================
% Section: Abstract  (outline.md "Abstract", 150-250 words)
% Draft v4 (2026-07-06). Wording mirrors the honest framing used in the body:
%   DPRNet matches CATANet-level accuracy with small gains on most benchmark cells
%   (1_introduction.tex, 3_experiments_comparison.tex, 3_experiments_efficiency.tex).
%   Numbers VERBATIM from paper/tables/* and paper/data/*: x4 Urban100 26.89 dB,
%   52.14G Multi-Adds (256x256-input protocol; below SwinIR-light 60.3G / SRFormer-light
%   56.5G, above CATANet 46.8G). Ablation is a 250k, 3-seed study: C1 dynamic
%   prototypes > EMA center, C4 balance loss -> more balanced + lower-variance usage.
%   No overstated compute/isolated-gain claims; consistent with method-experiment-traceability.md.
% =====================================================================
\begin{abstract}
Lightweight single image super-resolution (SR) must balance reconstruction quality with
model cost on resource-constrained platforms. Content-aware routing improves this balance
by grouping semantically related tokens, yet the widely used center-based router---the
Token Aggregation Block (TAB) of CATANet---relies on cross-batch exponential-moving-average
centers, hard assignment labels, and a coupled prototype update/confirmation step, tying
routing to a drifting historical state and discarding assignment confidence. We revisit
token routing itself and propose \emph{Dynamic Prototype Routing} (DPR), a confidence-aware,
drop-in replacement for the TAB routing that keeps the CATANet backbone and tensor interface
unchanged. DPR builds per-image prototypes directly
from the current tokens, decouples prototype generation from confirmation through a
learnable query bank, matches tokens with a bounded learnable temperature, and carries the
resulting assignment confidence into token ordering, global gating, and a soft fallback,
with a usage-balancing regularizer that discourages prototype-slot collapse. Trained on
DIV2K and evaluated on five standard benchmarks at $\times2/\times3/\times4$, DPRNet matches
CATANet-level accuracy with small gains on most benchmark cells while staying in the
lightweight regime; at $\times4$, for example, it obtains $26.89$\,dB on Urban100 with
$52.14$\,G Multi-Adds, fewer than the lightweight Transformer baselines SwinIR-light and
SRFormer-light under the same input convention. Three-seed ablations isolate each design
choice: the dynamic prototypes outperform a history-center EMA variant under an identical
protocol, and the balancing loss makes prototype usage more balanced and more stable across
seeds. DPR thus offers a reusable, interpretable alternative to center-based token routing
rather than a larger backbone.
\end{abstract}
